\documentclass[aspectratio=169,xcolor=table]{beamer}
\usepackage[utf8]{inputenc}
\usepackage[T1]{fontenc}
\usepackage{lmodern}
\usepackage{csquotes}
\usepackage{xcolor}
\usepackage[portuguese]{babel}
\usepackage{booktabs}
\usepackage{pifont}
\usepackage{tikz}
\usetikzlibrary{arrows.meta, positioning, fit, backgrounds, calc}
\usetheme{DCC}
\graphicspath{{imgs/}}

\newcommand{\cmark}{\ding{51}}
\newcommand{\xmark}{\ding{55}}

\author{Gustavo Santos \\ \vspace{0.2cm} {\small Orientador: Prof. Maycon L. M. Peixoto}}
\title{Desenvolvimento de um Serviço de Metadata Compatível com Cloud-Init para Incus}
\institute{Departamento de Ciência da Computação -- Universidade Federal da Bahia}
\date{\today}

\begin{document}

\begin{frame}[plain,noframenumbering]
    \titlepage
\end{frame}

\begin{frame}[plain,noframenumbering]{Agenda}
    \tableofcontents
\end{frame}

% =====================================================================
\section{Introdução}
\sectionframe

\begin{frame}{Contexto}
    \begin{itemize}
        \item Nuvens públicas (AWS, GCP, Azure) provisionam instâncias de forma
              programática e automatizada.
        \item Peça central: \textit{Instance Metadata Service} (IMDS), acessível pelo
              endereço \textit{link-local} \texttt{169.254.169.254}~\cite{aws_imds}.
        \item Cada instância descobre a própria configuração no \textit{boot}:
              hostname, endereços IP, chaves SSH, \textit{user-data}.
        \item O \textbf{cloud-init}~\cite{cloudinit2024} consome esses dados e é o
              padrão da indústria para inicialização automatizada de instâncias.
    \end{itemize}
\end{frame}

\begin{frame}{Problema}
    \begin{itemize}
        \item O \textbf{Incus}~\cite{incus2024} (\textit{fork} comunitário do LXD) não
              oferece um serviço de metadados HTTP nativo.
        \item A entrega de configuração existente (\textit{datasource} LXD via
              \textit{socket} \texttt{/dev/incus/sock}) sofre uma \textbf{condição de
              corrida}: o \textit{Incus agent} inicia tarde demais para os primeiros
              estágios do cloud-init em VMs.
        \item Consequências:
        \begin{itemize}
            \item configuração manual ou \textit{scripts} ad hoc;
            \item sem portabilidade de configurações entre nuvem pública e privada;
            \item OpenStack resolve, mas exige a plataforma completa — desproporcional
                  para \textit{homelabs}, borda e nuvens privadas pequenas.
        \end{itemize}
    \end{itemize}
\end{frame}

\begin{frame}{Objetivo e Contribuições}
    \textbf{Objetivo:} projetar e implementar um serviço de metadados compatível com
    cloud-init para o Incus, leve e autônomo.
    \vspace{0.4cm}
    \begin{enumerate}
        \item Análise das implementações de serviços de metadados em nuvens públicas e
              privadas e dos requisitos mínimos de compatibilidade com o cloud-init.
        \item Projeto e implementação do serviço: API REST, sincronização de
              instâncias orientada a eventos, persistência e alta disponibilidade.
        \item Avaliação em ambiente Incus real: corretude funcional, desempenho,
              escalabilidade e \textit{failover}.
    \end{enumerate}
\end{frame}

% =====================================================================
\section{Trabalhos Relacionados}
\sectionframe

\begin{frame}{Serviços de Metadados Existentes}
    \begin{itemize}
        \item \textbf{AWS IMDS}~\cite{aws_imds}: API HTTP em \texttt{169.254.169.254};
              IMDSv2 adiciona autenticação por sessão contra SSRF.
        \item \textbf{GCP Metadata Server}~\cite{gcp_metadata}: exige o cabeçalho
              \texttt{Metadata-Flavor: Google}; escopos de projeto e de instância.
        \item \textbf{OpenStack Metadata}~\cite{openstack_metadata}: compatível com a
              API EC2, mas fortemente acoplado ao Neutron e ao Nova — não executa de
              forma autônoma.
        \item \textbf{Datasource NoCloud}~\cite{cloudinit_nocloud}: obtém dados de uma
              URL base (\texttt{seedfrom}) consultando \texttt{/meta-data},
              \texttt{/user-data}, \texttt{/vendor-data} e \texttt{/network-config} —
              exige apenas \texttt{instance-id} e \texttt{local-hostname}.
        \item \textbf{RAFT}~\cite{ongaro2014raft}: consenso distribuído compreensível
              (eleição de líder + replicação de log); base da alta disponibilidade.
    \end{itemize}
\end{frame}

\begin{frame}{Comparação entre as Abordagens}
    \begin{table}
        \centering
        \small
        \renewcommand{\arraystretch}{1.2}
        \def\cellZ{\cellcolor{black!10}\xmark}
        \def\cellO{\cmark}
        \resizebox{\textwidth}{!}{%
        \begin{tabular}{lcccccc}
            \toprule
            \textbf{Solução} & \textbf{Autônomo} & \textbf{Compat. cloud-init} &
            \textbf{Contêineres} & \textbf{VMs} & \textbf{Alta Disp.} & \textbf{Cód. Aberto} \\
            \midrule
            AWS IMDS            & \cellZ & \cellO & \cellZ & \cellO & \cellO & \cellZ \\
            GCP Metadata        & \cellZ & \cellO & \cellZ & \cellO & \cellO & \cellZ \\
            OpenStack Metadata  & \cellZ & \cellO & \cellZ & \cellO & \cellO & \cellO \\
            LXD/Incus (atual)   & \cellO & \cellZ & \cellO & \cellO & \cellZ & \cellO \\
            \textbf{Este trabalho} & \cellO & \cellO & \cellO & \cellO & \cellO & \cellO \\
            \bottomrule
        \end{tabular}%
        }
    \end{table}
    \vspace{0.3cm}
    \textbf{Lacuna:} não há serviço de metadados leve e autônomo que implemente o
    protocolo HTTP esperado pelo cloud-init para o Incus.
\end{frame}

% =====================================================================
\section{Proposta}
\sectionframe

\begin{frame}{Visão Geral}
    \begin{itemize}
        \item Serviço de metadados compatível com cloud-init para instâncias Incus
              (contêineres e VMs).
        \item API REST consultada pelas instâncias durante o \textit{boot}, seguindo o
              modelo do \textit{datasource} NoCloud HTTP.
        \item A instância é identificada pelo \textbf{endereço IP de origem} da
              requisição — sem agentes ou configuração dentro da instância.
        \item Sincronização orientada a eventos mantém o estado consistente com o
              hipervisor Incus.
        \item Implantação autônoma: um único binário Go, banco embutido, sem
              dependências externas.
    \end{itemize}
\end{frame}

\begin{frame}{Arquitetura}
    \begin{figure}
        \centering
        \resizebox{!}{0.72\textheight}{%
        \begin{tikzpicture}[
            node distance=1.6cm and 2.2cm,
            every node/.style={font=\small},
            box/.style={draw, rounded corners=3pt, minimum width=3cm, minimum height=1cm, align=center, fill=white, thick},
            extbox/.style={draw, rounded corners=3pt, minimum width=3cm, minimum height=1cm, align=center, fill=black!8, thick},
            groupbox/.style={draw, dashed, rounded corners=6pt, inner sep=14pt, thick, fill=white},
            arr/.style={-{Stealth[length=6pt]}, thick},
        ]
            \node[extbox] (incus) {Incus API};
            \node[extbox, right=7cm of incus] (instances) {Instâncias\\(VMs / Contêineres)};

            \node[box, below=1.8cm of $(incus)!0.5!(instances)$] (server) {\textbf{cmd/server}\\Inicialização};

            \node[box, below left=1.6cm and 2.0cm of server] (events) {\textbf{events}\\Cron + EventBus};
            \node[box, below right=1.6cm and 2.0cm of server] (api) {\textbf{api}\\Gin (REST)};

            \node[box, below=1.8cm of $(events)!0.5!(api)$] (storage) {\textbf{storage}\\SQLite + SQLC};
            \node[box, left=2.0cm of storage] (consensus) {\textbf{consensus}\\RAFT};
            \node[box, right=2.0cm of storage] (telemetry) {\textbf{telemetry}\\OTel + Prometheus};

            \begin{scope}[on background layer]
                \node[groupbox, fill=blue!3,
                      fit=(server)(events)(api)(storage)(consensus)(telemetry),
                      label={[font=\footnotesize\bfseries, anchor=north west, yshift=-3pt]north west:Serviço de Metadados}] (service) {};
            \end{scope}

            \draw[arr, gray] (server) -- (events);
            \draw[arr, gray] (server) -- (api);
            \draw[arr] (incus.south) -- node[right, font=\scriptsize, pos=0.35] {consulta instâncias} ++(0,-1.0) -| (events.north);
            \draw[arr] (instances.south) -- node[left, font=\scriptsize, pos=0.35] {GET /configs/*} ++(0,-1.0) -| (api.north);
            \draw[arr] (events.south east) -- node[right, font=\scriptsize, pos=0.4] {escrita} (storage.north west);
            \draw[arr] (api.south west) -- node[left, font=\scriptsize, pos=0.4] {leitura} (storage.north east);
            \draw[arr] (events.south west) -- node[left, font=\scriptsize, pos=0.4] {replica} (consensus.north);
            \draw[arr] (consensus.east) -- node[below, font=\scriptsize] {aplica} (storage.west);
            \draw[arr, gray, dashed] (events.south east) -- (telemetry.north west);
            \draw[arr, gray, dashed] (api.south east) -- node[right, font=\scriptsize, pos=0.3, text=gray] {métricas} (telemetry.north);
        \end{tikzpicture}%
        }
    \end{figure}
\end{frame}

\begin{frame}{Sincronização de Instâncias}
    Ciclo periódico e orientado a eventos (\textit{pub/sub}):
    \begin{enumerate}
        \item \textit{Cron} a cada 10~s publica o evento \texttt{instances.sync}.
        \item O \textit{handler} consulta a API do Incus e lista as instâncias ativas.
        \item Para cada instância, um evento \texttt{instance.sync} é publicado —
              processamento desacoplado e paralelo.
        \item Para cada instância, o serviço:
        \begin{itemize}
            \item extrai rede (IPv4, IPv6, MAC) do estado reportado pelo Incus;
            \item persiste metadados cloud-init (\texttt{instance-id}, hostname,
                  chaves SSH, \textit{placement});
            \item armazena \texttt{user-data} e \texttt{vendor-data} das chaves
                  \texttt{cloud-init.*} do Incus;
            \item gera \texttt{network-config} (Networking Config v2 / Netplan) quando
                  não definida pelo administrador.
        \end{itemize}
    \end{enumerate}
\end{frame}

\begin{frame}{Endpoints da API}
    Endpoints públicos, no modelo do \textit{datasource} NoCloud
    HTTP~\cite{cloudinit_nocloud}:
    \begin{itemize}
        \item \texttt{GET /configs/meta-data} — identidade da instância
              (\texttt{instance-id}, hostname, IPs, \texttt{public-keys}).
        \item \texttt{GET /configs/meta-data/:key} — campo específico dos metadados.
        \item \texttt{GET /configs/user-data} — \textit{scripts}/cloud-config do usuário.
        \item \texttt{GET /configs/vendor-data} — por instância, com \textit{fallback}
              global.
        \item \texttt{GET /configs/network-config} — definida pelo administrador ou
              gerada automaticamente.
    \end{itemize}
    \vspace{0.3cm}
    Endpoints internos (\texttt{/internal/*}): CRUD administrativo de vendor-data —
    não expostos às instâncias.
\end{frame}

\begin{frame}{Descoberta do Serviço pelas Instâncias}
    \begin{itemize}
        \item Mesma convenção das nuvens públicas: endereço \textit{link-local}
              \texttt{169.254.169.254} adicionado à \textit{bridge} do Incus
              (\texttt{incusbr0}).
        \item Regra NAT redireciona \texttt{169.254.169.254:80} para a porta do
              serviço — nenhuma configuração dentro das instâncias.
        \item \textit{Datasource} NoCloud em modo rede (\texttt{dsmode: net}),
              com \texttt{seedfrom} apontando para \texttt{/configs/} no perfil do
              Incus.
    \end{itemize}
\end{frame}

\begin{frame}{Alta Disponibilidade via RAFT}
    Nó único = ponto único de falha. Solução: consenso distribuído
    RAFT~\cite{ongaro2014raft} com \texttt{hashicorp/raft}~\cite{hashicorp_raft}.
    \begin{itemize}
        \item \textbf{Eleição de líder}: apenas o líder sincroniza com o Incus e propõe
              escritas; falha do líder dispara nova eleição automática.
        \item \textbf{Replicação de log + FSM}: cada escrita vira um comando tipado,
              confirmado por maioria e aplicado ao SQLite local de cada nó.
        \item \textbf{Persistência}: log RAFT em BoltDB — sem banco externo.
        \item \textbf{Leituras em qualquer nó}: todos mantêm cópia local do estado,
              distribuindo a carga de consultas.
        \item Ativação opcional por variáveis de ambiente
              (\texttt{RAFT\_ENABLED}, \texttt{NODE\_ID}, \texttt{PEERS}, \ldots).
    \end{itemize}
\end{frame}

\begin{frame}{Tecnologias Utilizadas}
    \begin{itemize}
        \item \textbf{Go} — binário estático único; goroutines para o processamento
              paralelo de eventos.
        \item \textbf{Gin}~\cite{gin2024} — framework HTTP de alto desempenho; separa
              rotas públicas e internas.
        \item \textbf{SQLite + SQLC}~\cite{sqlite2024} — banco embutido; queries SQL
              com tipagem estática gerada.
        \item \textbf{GoEventBus} — barramento \textit{pub/sub} que desacopla
              agendamento e sincronização.
        \item \textbf{gocron} — agendamento do ciclo de sincronização (10~s).
        \item \textbf{hashicorp/raft}~\cite{hashicorp_raft} — consenso distribuído.
        \item \textbf{OpenTelemetry + Prometheus}~\cite{opentelemetry2024,prometheus2024}
              — métricas de sincronização e de API em \texttt{/metrics}.
    \end{itemize}
\end{frame}

% =====================================================================
\section{Resultados}
\sectionframe

\begin{frame}{Metodologia de Avaliação}
    \begin{itemize}
        \item Ambiente: GCE \texttt{e2-standard-4} (4 vCPUs, 15~GiB), Ubuntu 24.04
              LTS, Incus 6.0.0; serviço no mesmo host das instâncias.
        \item Cargas de 5 a 200 contêineres simultâneos com cloud-init apontando para
              o serviço.
        \item Latência medida com \texttt{hey}: 5000 requisições, concorrência 50;
              percentis p50/p95/p99.
        \item \textbf{Achado prático}: em nuvem pública, \texttt{169.254.169.254} já é
              usado pelo provedor — o \textit{datasource} foi apontado para o
              \textit{gateway} da \textit{bridge} (\texttt{10.10.10.1:8080}).
    \end{itemize}
\end{frame}

\begin{frame}{Avaliação Funcional}
    \begin{table}
        \centering
        \begin{tabular}{lc}
            \toprule
            \textbf{Propriedade} & \textbf{Resultado} \\
            \midrule
            Isolamento de metadados & Aprovado \\
            Isolamento sob falsificação de IP (\texttt{X-Forwarded-For}) & Aprovado \\
            Execução de \textit{user-data} & Aprovado \\
            Entrega de \textit{network-config} & Aprovado \\
            Propagação de estado após reinício & Aprovado \\
            \bottomrule
        \end{tabular}
    \end{table}
    \vspace{0.3cm}
    Inicialização real: cloud-init selecionou o \textit{datasource} NoCloud de rede,
    obteve os dados (HTTP 200) e concluiu com \texttt{cloud-init status: done} —
    compatibilidade de ponta a ponta.
\end{frame}

\begin{frame}{Desempenho: Latência da API}
    \begin{table}
        \centering
        \begin{tabular}{lccc}
            \toprule
            \textbf{Endpoint} & \textbf{p50 (ms)} & \textbf{p95 (ms)} & \textbf{p99 (ms)} \\
            \midrule
            \texttt{/meta-data}       & 13{,}2 & 53{,}7 & 83{,}2 \\
            \texttt{/user-data}       & 18{,}9 & 53{,}6 & 72{,}4 \\
            \texttt{/network-config}  & 11{,}4 & 47{,}2 & 72{,}4 \\
            \texttt{/vendor-data}     & 17{,}0 & 48{,}4 & 68{,}9 \\
            \bottomrule
        \end{tabular}
    \end{table}
    \vspace{0.3cm}
    \begin{itemize}
        \item Medianas de 11--19~ms; p99 abaixo de 90~ms — compatível com as janelas
              de tempo do \textit{boot}.
        \item Latência de sincronização (criação no Incus $\rightarrow$ metadados
              disponíveis): $\approx$~20{,}5~s ($\approx$~2 ciclos de \textit{polling}).
    \end{itemize}
\end{frame}

\begin{frame}{Desempenho: Escalabilidade}
    \begin{table}
        \centering
        \small
        \begin{tabular}{lccccc}
            \toprule
            \textbf{Instâncias} & \textbf{p50 (ms)} & \textbf{p95 (ms)} & \textbf{p99 (ms)} & \textbf{Erros} & \textbf{Mem. (MB)} \\
            \midrule
            5   & 16{,}2 & 60{,}3 & 92{,}4 & 0{,}0\% & 1152 \\
            25  & 12{,}4 & 51{,}8 & 77{,}2 & 0{,}0\% & 1373 \\
            50  & 12{,}9 & 52{,}0 & 80{,}8 & 0{,}0\% & 1733 \\
            100 & 14{,}4 & 58{,}6 & 92{,}1 & 0{,}0\% & 1763 \\
            200 & 12{,}8 & 54{,}2 & 87{,}2 & 0{,}0\% & 1711 \\
            \bottomrule
        \end{tabular}
    \end{table}
    \vspace{0.3cm}
    \begin{itemize}
        \item Hipótese inicial: degradação por serialização de escrita do SQLite.
        \item \textbf{Não corroborada} até 200 instâncias: latência estável e zero
              erros — carga dominada por leituras; escritas a cada 10~s absorvidas
              pelo modo WAL.
    \end{itemize}
\end{frame}

\begin{frame}{Alta Disponibilidade (\textit{cluster} de 3 nós)}
    \begin{table}
        \centering
        \begin{tabular}{ll}
            \toprule
            \textbf{Propriedade} & \textbf{Resultado} \\
            \midrule
            Formação do \textit{cluster} & 1 líder + 2 seguidores \\
            Replicação de estado & instância presente nos 3 nós \\
            Reeleição após falha do líder & \textbf{2{,}40~s} \\
            Disponibilidade após \textit{failover} & preservada \\
            Reintegração de nó & reingressa como seguidor \\
            \bottomrule
        \end{tabular}
    \end{table}
    \vspace{0.3cm}
    \begin{itemize}
        \item Topologia avaliada: réplicas sobre um Incus compartilhado — reconciliação
              correta e \textit{failover} sem perda de dados.
        \item Um Incus por nó (borda distribuída) exigiria reconciliação por nó de
              origem.
    \end{itemize}
\end{frame}

\begin{frame}{Discussão e Limitações}
    \begin{itemize}
        \item Cobre os campos essenciais dos módulos mais comuns do cloud-init
              (\texttt{set\_hostname}, \texttt{users-groups}, \texttt{write-files},
              \texttt{runcmd}, rede).
        \item \textbf{Limitações}:
        \begin{itemize}
            \item SQLite serializa escritas — potencial gargalo bem acima da escala
                  avaliada;
            \item subconjunto dos campos de AWS/GCP (sem tipo de instância,
                  faturamento);
            \item garantias de \textit{failover} pressupõem a topologia de Incus
                  compartilhado.
        \end{itemize}
        \item \textbf{Ameaças à validade}: host único nos testes de desempenho;
              contêineres de sistema (não VMs completas); janela de inconsistência do
              \textit{polling}.
    \end{itemize}
\end{frame}

% =====================================================================
\section{Conclusão}
\sectionframe

\begin{frame}{Conclusão}
    \begin{itemize}
        \item Serviço de metadados compatível com cloud-init para o Incus: preenche a
              lacuna de um endpoint HTTP autônomo em infraestrutura privada.
        \item Identificação por IP de origem: operação transparente, sem agentes nas
              instâncias.
        \item Validado em ambiente real: \textit{boots} completos com sucesso,
              latências de dezenas de milissegundos, escalabilidade estável até 200
              instâncias e reeleição de líder em $\approx$~2{,}4~s.
        \item Código aberto, binário único, implantação simples — adequado a
              \textit{homelabs}, borda e nuvens privadas de pequena escala.
    \end{itemize}
\end{frame}

\begin{frame}{Trabalhos Futuros}
    \begin{enumerate}
        \item Suporte à convenção \texttt{169.254.169.254} via \textit{namespace} de
              rede ou \textit{proxy} reverso.
        \item Ampliação dos módulos e campos do cloud-init suportados.
        \item Autenticação baseada em sessões inspirada no
              IMDSv2~\cite{aws_imdsv2} contra SSRF.
        \item Substituição do \textit{polling} pela API de eventos do Incus
              (sincronização em tempo real).
        \item \textit{Benchmarks} em cenários de maior escala.
    \end{enumerate}
\end{frame}

\begin{frame}[plain,noframenumbering]
    \centering
    \vspace{2cm}
    {\LARGE \textbf{Obrigado!}}\\[1cm]
    {\large Perguntas?}\\[1.5cm]
    \texttt{gustavoafs@ufba.br}
\end{frame}

\begin{frame}[plain,noframenumbering,allowframebreaks]
    \frametitle{Referências}
    \bibliographystyle{plain}
    \bibliography{refs}
\end{frame}

\end{document}
